% page 426 du fac-simile (image p-425 du scan)
% bloc 160.0 x 250.6 mm ; marge gauche 8.0 mm
\pgc{-12.700mm}{0.000mm}{
\l{}
\l{}
\l{}
\l{\cel{3}418.}
\l{}
\l{\sou{palpebro}. (\sou{anat.})\cel{2}Singla ek la du membrani movebla qui,}
\l{\cel{3}per interkontaktar, tegas la globo okulala. - FIS.}
\l{}
\l{\sou{palpitar}. (\sou{netrans}.) (\sou{aludante parto di la homo-korpo}).}
\l{\cel{3}Anhelar konvulsatre. Pulsar konvulsatre. - DEFIS.}
\l{}
\l{\sou{paltoto}. Vesto quan la viri +weras super frako o redingoto.}
\l{\cel{3}DEFRS.}
\l{}
\l{\sou{palumbo}. Kolombo sovaja, qua lojas en nesto sur la arbori.}
\l{\cel{3}- FIS.}
\l{}
\l{\sou{pamfleto}. Verketo violentoza en qua onu atakas ulu, ulo.}
\l{\cel{3}DFR.}
\l{}
\l{\sou{pampro}. La foliaro di vito.}
\l{}
\l{\sou{pano}. Nutrivo ek maso de farino dilutita en aquo, petrisita}
\l{\cel{3}e bakita. - FIS.}
\l{}
\l{\sou{panaceo}. (\sou{medicino}). Remediivo generala, unika, apta por}
\l{\cel{3}irga morbo. - DEFIRS.}
\l{}
\l{\sou{panariso}. (\sou{patol.)}\cel{2}Inflamuro flegmona ye la extremajo di}
\l{\cel{3}la fingri. - EFIS.}
\l{}
\l{\sou{pancho}.\cel{2}La maxim kapaca ek la stomaki di la rumineri. -DEFIS}
\l{}
\l{\sou{pandur}o. Soldato exter-armeana, privata. - DEFIRS.}
\l{}
\l{\sou{panear}. (\sou{netrans}.)Ne plus povar avancar, funcionar. - DFI.}
\l{}
\l{\sou{panegir}o.Diskurso publika laude di persono. - DEFIRS.}
\l{}
\l{\sou{panelo.}\cel{1}I.(arkitekt.)\cel{2}Parto plu o min granda e grosa di muro.-}
\l{\cel{3}Ensemblo de karpenturo di qua onu plenigas la vakui ye}
\l{\cel{3}masonuro. - II. Peco ek vesto, ek mantelo, ek robo, e c.}
\l{\cel{3}DEFR.}
\l{}
\l{\sou{pangenezo(biol.)}\cel{3}Teorio per qua Darwin explikas omna feno-}
\l{\cel{3}meni biologiala, per la karakterizivi e per la migri da}
\l{\cel{3}partikuli materia, quin lu nomizas \textquotedbl{}jemuli\textquotedbl{}. -DEFIS.}
\l{}
\l{\sou{pangeno.}\cel{3}(\sou{biol.)}\cel{2}Segun la teorio di de Vries : mikra kor-}
\l{\cel{3}po, ek molekulo kemiala, ma dotita ye la karakterizivo}
\l{\cel{3}di la materio vivanta : asimiliveso, kreskiveso, su-}
\l{\cel{3}reprodukto per su-divido. - DEFIS.}
\l{}
\l{\sou{pangolino.}\cel{3}(\sou{zool.})\cel{2}Sen-dento, di qua la korpo esassqua-}
\l{\cel{3}moza. ek Sud-Afrika ed ek India. - L.\cel{1}\sou{manis laticaudata}.}
\l{}
\l{\sou{panio}.- Peco de stofo per qua la negri mi-nuda tegas su}
\l{\cel{3}de la talio til la genui. - F.}
}
